\section{The rigidity lemma and its corollaries}
\label{mk-chap:AbelianVarietyRigidity}
\label{mk-chap:avr_for_rr}

This chapter proves the rigidity lemma and the two elementary consequences used in the
Albanese construction.  The argument is characteristic-free and uses only properness,
affine constancy, and separatedness.  In particular, it does not use the theorem of the
cube or coherent cohomology.

Throughout, \(\bar k\) is an algebraically closed field.  Products are fibre products over
\(\Spec \bar k\), and a point means a \(\bar k\)-point.
Following the rigidity arguments of Mumford and Milne, we use the formulation in
\cite{mumford-av,milne-av}.

\section{The rigidity lemma}

\begin{lemma}[Rigidity Lemma (Mumford, Form I)]
  \label{mk-thm:rigidity_lemma}
  \lean{AlgebraicGeometry.rigidity_lemma}
  \leanok
  \source{mumford-abelian-varieties:page-0054}
  \dcref{custom}

  Let \(X\) be a complete variety, let \(Y\) and \(Z\) be varieties, fix a point
  \(x_0\in X\), and let \(f \colon X \times Y \to Z\) be a morphism.  Suppose that
  for some \(y_0 \in Y\) the image \(f(X \times \{y_0\})\) is the single point
  \(z_0\).  Then there is a morphism \(g \colon Y \to Z\) such that
  \[
    f = g \circ p_2,
  \]
  where \(p_2 \colon X \times Y \to Y\) is the second projection.
\end{lemma}
\begin{proof}
  \leanok
  \uses{mk-lem:rigidity_eqOn_dense_open, mk-lem:rigidity_core, mk-lem:rigidity_snd_lift}
  Define \(g(y)=f(x_0,y)\).  Let \(U_0\) be an affine
  neighbourhood of \(z_0\), put \(F=Z\setminus U_0\), and set
  \[
    G=p_2\bigl(f^{-1}(F)\bigr).
  \]
  The projection \(p_2\) is closed because \(X\) is complete, hence \(G\) is closed.
  The hypothesis on the fibre over \(y_0\) gives \(y_0\notin G\), so
  \(V=Y\setminus G\) is nonempty.  For every \(y\in V\), the complete variety
  \(X\times\{y\}\) maps into the affine variety \(U_0\), and therefore its image is
  a point.  That point is \(f(x_0,y)\).  Thus \(f\) and \(g\circ p_2\) agree on the
  nonempty open set \(X\times V\).  This open set is dense, and the target is
  separated, so the two morphisms agree everywhere.
\end{proof}

\begin{remark}
  \label{mk-rmk:rigidity_lemma_cube_free}
  \dcref{custom}
  The collapse assumption is essential.  For example, the first projection
  \(X\times Y\to X\) does not agree on a nonempty open set with the map obtained by
  replacing the first coordinate by a fixed point, unless \(X\) is itself a point.
  The proof above uses neither the theorem of the cube nor cohomology.
\end{remark}

\begin{remark}
  \label{mk-rmk:rigidity_lemma_decomposition}
  \dcref{custom}
  The proof separates into four ingredients: the cartesian identity for the collapsed
  map, closedness of the second projection, constancy of a proper integral slice mapping
  to an affine scheme, and extension of equality from a dense set to a separated target.
  The following lemmas record these ingredients individually.
\end{remark}

\begin{lemma}[Cartesian collapse identity]
  \label{mk-lem:rigidity_snd_lift}
  \lean{AlgebraicGeometry.rigidity_snd_lift}
  \leanok
  \dcref{custom}

  Fix \(x_0\in X\), and let \(s_{x_0}\colon Y\to X\times Y\) be the section
  \(y\mapsto(x_0,y)\) of \(p_2\).  Then \(s_{x_0}\circ p_2\) is the endomorphism
  \((x,y)\mapsto(x_0,y)\) of \(X\times Y\).
\end{lemma}
\begin{proof}
  \leanok
  The second component is the identity because \(p_2\circ s_{x_0}=\mathrm{id}_Y\).
  The first component is the constant map with value \(x_0\).  The product universal
  property therefore identifies the composite with \((x,y)\mapsto(x_0,y)\).
\end{proof}

\begin{lemma}[Completeness makes the second projection closed]
  \label{mk-lem:snd_left_isClosedMap}
  \lean{AlgebraicGeometry.snd_left_isClosedMap}
  \leanok
  \dcref{custom}

  If \(X\) is proper over \(\bar k\), then the underlying topological map of
  \(p_2\colon X\times Y\to Y\) is closed.
\end{lemma}
\begin{proof}
  \leanok
  The projection is the base change of \(X\to\Spec\bar k\) along
  \(Y\to\Spec\bar k\).  Proper morphisms are universally closed, and universal
  closedness is stable under base change.  Hence \(p_2\) is closed.
\end{proof}

\begin{lemma}[Scheme-theoretic gluing core]
  \label{mk-lem:rigidity_core}
  \lean{AlgebraicGeometry.rigidity_core}
  \leanok
  \dcref{custom}

  Suppose that \(X\) is proper, \(X\times Y\) is reduced, geometrically irreducible,
  and locally of finite type over \(\bar k\), and \(Z\) is separated.  Fix points
  \(x_0\in X\), \(y_0\in Y\), and \(z_0\in Z\).  If
  \(f\colon X\times Y\to Z\) maps \(X\times\{y_0\}\) to \(z_0\), then \(f\)
  equals the morphism \((x,y)\mapsto f(x_0,y)\) on all of \(X\times Y\).
\end{lemma}
\begin{proof}
  \leanok
  \uses{mk-lem:rigidity_eqOn_dense_open}
  By \ref{mk-lem:rigidity_eqOn_dense_open}, the collapse hypothesis gives a nonempty
  open set on which the two morphisms agree.  Geometric irreducibility makes this open
  set dense.  Two morphisms from a reduced scheme to a separated scheme that agree
  after a dominant open immersion are equal, so the equality holds everywhere.
\end{proof}

\begin{lemma}[Dense-open agreement]
  \label{mk-lem:rigidity_eqOn_dense_open}
  \lean{AlgebraicGeometry.rigidity_eqOn_dense_open}
  \leanok
  \source{mumford-abelian-varieties:page-0054}
  \dcref{custom}

  Let \(X\) be proper, suppose that \(X\times Y\) is reduced, geometrically
  irreducible, and locally of finite type over \(\bar k\), and let \(Z\) be separated.
  Fix points \(x_0\in X\), \(y_0\in Y\), and \(z_0\in Z\).  If
  \(f\colon X\times Y\to Z\) maps \(X\times\{y_0\}\) to \(z_0\), then there is a
  nonempty open subset on which
  \[
    f(x,y)=f(x_0,y).
  \]
\end{lemma}
\begin{proof}
  \leanok
  \uses{mk-lem:rigidity_eqOn_saturated_open_to_affine, mk-lem:snd_left_isClosedMap}
  Choose an affine open neighbourhood \(U_0\) of \(z_0\), set
  \(F=Z\setminus U_0\), and let \(G=p_2(f^{-1}(F))\).  By
  \ref{mk-lem:snd_left_isClosedMap}, the set \(G\) is closed.  The collapsed fibre is
  disjoint from \(f^{-1}(F)\), so \(y_0\notin G\).  Hence
  \(V=Y\setminus G\) is nonempty and \(U=p_2^{-1}(V)\) is a nonempty saturated
  open subset of \(X\times Y\) with \(f(U)\subseteq U_0\).  Apply
  \ref{mk-lem:rigidity_eqOn_saturated_open_to_affine} to \(U\).
\end{proof}

\begin{lemma}[Constancy on a saturated open]
  \label{mk-lem:rigidity_eqOn_saturated_open_to_affine}
  \lean{AlgebraicGeometry.rigidity_eqOn_saturated_open_to_affine}
  \leanok
  \source{mumford-abelian-varieties:page-0054}
  \dcref{custom}

  Let \(X\) be proper, suppose that \(X\times Y\) is reduced, geometrically
  irreducible, and locally of finite type over \(\bar k\), and let \(Z\) be separated.
  Fix \(x_0\in X\), and let \(U=p_2^{-1}(V)\) be a saturated open subset of
  \(X\times Y\).  If \(f\colon X\times Y\to Z\) maps \(U\) into an affine open
  subset \(U_0\subseteq Z\), then on \(U\) one has \(f(x,y)=f(x_0,y)\).
\end{lemma}
\begin{proof}
  \leanok
  \uses{mk-lem:rigidity_eqAt_closedPoint_of_proper_into_affine,
    mk-lem:morphism_eq_of_eqAt_closedPoints}
  The open subscheme \(U\) is locally of finite type over the Jacobson scheme
  \(\Spec\bar k\), hence its closed points are dense.  At every closed point of
  \(U\), the desired equality follows from
  \ref{mk-lem:rigidity_eqAt_closedPoint_of_proper_into_affine}.  The source is reduced
  and the target is separated, so
  \ref{mk-lem:morphism_eq_of_eqAt_closedPoints} promotes these pointwise equalities to
  equality of morphisms on \(U\).
\end{proof}

\begin{lemma}[Equality from closed points]
  \label{mk-lem:morphism_eq_of_eqAt_closedPoints}
  \lean{AlgebraicGeometry.morphism_eq_of_eqAt_closedPoints}
  \leanok
  \dcref{custom}

  Let \(W\) be a reduced Jacobson scheme and let \(Z\) be separated.  If two
  morphisms \(g_1,g_2\colon W\to Z\) agree after restriction along
  \(\Spec\kappa(x)\to W\) for every closed point \(x\in W\), then \(g_1=g_2\).
\end{lemma}
\begin{proof}
  \leanok
  Form the coproduct
  \[
    P=\coprod_{x\in W_{\mathrm{cl}}}\Spec\kappa(x)
  \]
  and its natural morphism \(q\colon P\to W\).  Its image contains every closed
  point of \(W\), so it is dense because \(W\) is Jacobson.  Thus \(q\) is dominant.
  The hypotheses and the coproduct universal property give
  \(g_1\circ q=g_2\circ q\).  Dominance of \(q\), reducedness of \(W\), and
  separatedness of \(Z\) then imply \(g_1=g_2\).
\end{proof}

\begin{lemma}[Proper integral schemes are constant on affine targets]
  \label{mk-lem:eq_comp_of_isAffine_of_properIntegral}
  \lean{AlgebraicGeometry.eq_comp_of_isAffine_of_properIntegral}
  \leanok
  \dcref{custom}

  Let \(W\) be an integral scheme, proper and locally of finite type over \(\bar k\),
  and let \(V\) be affine.  If \(a,b\colon\Spec\bar k\to W\) are sections of the
  structure morphism and \(g\colon W\to V\), then \(g\circ a=g\circ b\).
\end{lemma}
\begin{proof}
  \leanok
  The ring \(\Gamma(W,\mathcal O_W)\) is a field because \(W\) is integral and
  proper.  It is finite over \(\bar k\) because \(W\) is locally of finite type.
  Since \(\bar k\) is algebraically closed, the structure map
  \(\bar k\to\Gamma(W,\mathcal O_W)\) is an isomorphism.  The two sections induce
  the same inverse on global sections.  Morphisms into the affine scheme \(V\) are
  determined by their maps on global sections, so \(g\circ a=g\circ b\).
\end{proof}

\begin{lemma}[Integrality descends to a retract]
  \label{mk-lem:isIntegral_of_retract_of_integral}
  \lean{AlgebraicGeometry.isIntegral_of_retract}
  \leanok
  \dcref{custom}

  Let \(T\) be an integral scheme.  If morphisms \(r\colon S\to T\) and
  \(p\colon T\to S\) satisfy \(p\circ r=\mathrm{id}_S\), then \(S\) is integral.
\end{lemma}
\begin{proof}
  \leanok
  The map \(p\) is surjective because it has the section \(r\).  Therefore the
  continuous image of the irreducible space \(T\) is all of \(S\), and \(S\) is
  irreducible.  For each \(x\in S\), functoriality of stalk maps and
  \(p\circ r=\mathrm{id}_S\) show that
  \[
    \mathcal O_{S,x}\longrightarrow\mathcal O_{T,r(x)}
  \]
  is split injective.  The target is reduced, so \(\mathcal O_{S,x}\) is reduced.
  Hence every stalk of \(S\) is reduced, and \(S\) is reduced.  Reducedness and
  irreducibility give integrality.
\end{proof}

\begin{lemma}[Constancy on a closed slice]
  \label{mk-lem:rigidity_eqAt_closedPoint_of_proper_into_affine}
  \lean{AlgebraicGeometry.rigidity_eqAt_closedPoint_of_proper_into_affine}
  \leanok
  \source{mumford-abelian-varieties:page-0054}
  \dcref{custom}

  In the setting of
  \ref{mk-lem:rigidity_eqOn_saturated_open_to_affine}, let \(x\) be a closed point of
  \(U\).  Then \(f\) and the map \((x',y)\mapsto f(x_0,y)\) agree after restriction
  along the residue-field point \(\Spec\kappa(x)\to U\).
\end{lemma}
\begin{proof}
  \leanok
  \uses{mk-lem:eq_comp_of_isAffine_of_properIntegral,
    mk-lem:isIntegral_of_retract_of_integral}
  Since \(U\) is locally of finite type over the algebraically closed field \(\bar k\),
  the residue field of its closed point \(x\) is \(\bar k\).  Thus \(x\) is a
  \(\bar k\)-point.  Write \(y=p_2(x)\); then \(y\) and \((x_0,y)\) are also
  \(\bar k\)-points.  Since \(U\) is saturated, the whole fibre \(X_y\) lies in
  \(U\), and its image under \(f\) lies in the affine open \(U_0\).  Identifying
  \(X_y\) with \(X\), the slice section and first projection exhibit \(X\) as a
  retract of the integral product \(X\times Y\), so it is integral by
  \ref{mk-lem:isIntegral_of_retract_of_integral}; it is also proper and locally of
  finite type over \(\bar k\).  The two points of \(X_y\) represented by \(x\) and
  \((x_0,y)\) therefore have the same image in \(U_0\) by
  \ref{mk-lem:eq_comp_of_isAffine_of_properIntegral}.  Restricting this equality along
  \(\Spec\kappa(x)\to U\) gives the claim.
\end{proof}

\section{Product decompositions and pointed maps}

\begin{remark}
  \label{mk-rmk:cube_not_needed}
  \dcref{custom}
  The results in this section are direct consequences of the rigidity lemma.  The theorem
  of the cube enters the theory of abelian varieties later and is not needed here.
\end{remark}

\begin{lemma}[Canonical decomposition over a product]
  \label{mk-lem:hom_additivity_over_product}
  \lean{AlgebraicGeometry.hom_additive_decomp_of_rigidity}
  \leanok
  \source{abelian-varieties:page-0015}
  \dcref{custom}

  Let \(V\) be proper, suppose that \(V\times W\) is a reduced geometrically
  irreducible variety, and let \(A\) be a proper group variety.  Fix points
  \(v_0\in V\), \(w_0\in W\).  If \(h\colon V\times W\to A\) satisfies
  \(h(v_0,w_0)=e\), define
  \[
    f(v)=h(v,w_0),\qquad g(w)=h(v_0,w).
  \]
  Then
  \[
    h(v,w)=f(v)g(w).
  \]
  Equivalently, \(h=(f\circ p)(g\circ q)\), where \(p\) and \(q\) are the two
  projections.  Commutativity of \(A\) is not required.
\end{lemma}
\begin{proof}
  \leanok
  \uses{mk-thm:rigidity_lemma}
  Set \(\varphi=h/((f\circ p)(g\circ q))\).  On \(V\times\{w_0\}\), the map
  \(\varphi\) is the identity element because \(g(w_0)=h(v_0,w_0)=e\).  By the
  rigidity lemma, \(\varphi\) factors through \(q\).  On the section
  \(\{v_0\}\times W\), it is again the identity, so the factor through \(q\) is
  identically the identity.  Hence \(\varphi=e\) and the displayed decomposition
  follows.
\end{proof}

\begin{corollary}[Unique additive decomposition (Milne Corollary 1.5)]
  \label{mk-cor:hom_additivity_over_product_unique}
  \dcref{1.5,custom}
  Let \(V\) and \(W\) be complete varieties with points \(v_0\) and \(w_0\), and let
  \(A\) be an abelian variety.  Every morphism \(h\colon V\times W\to A\) with
  \(h(v_0,w_0)=0\) has a unique expression
  \[
    h=f\circ p+g\circ q,
  \]
  where \(f(v_0)=0\) and \(g(w_0)=0\).
\end{corollary}
\begin{proof}
  \uses{mk-lem:hom_additivity_over_product}
  Existence is \ref{mk-lem:hom_additivity_over_product}.  Restricting a decomposition to
  \(V\times\{w_0\}\) recovers \(f\), and restricting it to
  \(\{v_0\}\times W\) recovers \(g\).  Thus the two summands are unique.
\end{proof}

\begin{lemma}[A pointed regular map is a homomorphism]
  \label{mk-lem:av_regular_map_is_hom}
  \lean{AlgebraicGeometry.av_regularMap_isHom_of_zero}
  \leanok
  \source{abelian-varieties:page-0015}
  \dcref{custom}

  Let \(A\) and \(B\) be abelian varieties, and let \(\alpha\colon A\to B\) be a
  regular map with \(\alpha(0_A)=0_B\).  Then \(\alpha\) is a group homomorphism.
\end{lemma}
\begin{proof}
  \leanok
  \uses{mk-lem:hom_additivity_over_product}
  Consider
  \[
    \varphi(a,a')=\alpha(a+a')-\alpha(a)-\alpha(a').
  \]
  This is a regular map \(A\times A\to B\), and it vanishes on both coordinate
  axes.  Apply \ref{mk-lem:hom_additivity_over_product} to \(\varphi\).  Both canonical
  summands are zero, hence \(\varphi=0\).  Therefore
  \(\alpha(a+a')=\alpha(a)+\alpha(a')\) for all \(a,a'\), so \(\alpha\) is a
  homomorphism.
\end{proof}

\begin{corollary}[Regular maps are translations of homomorphisms (Milne Corollary 1.2)]
  \label{mk-cor:av_regular_map_translation}
  \dcref{1.2,custom}
  Every regular map \(\alpha\colon A\to B\) of abelian varieties is the composite of
  a homomorphism with a translation.
\end{corollary}
\begin{proof}
  \uses{mk-lem:av_regular_map_is_hom}
  Put \(b=\alpha(0_A)\) and define \(\beta(a)=\alpha(a)-b\).  Then
  \(\beta(0_A)=0_B\), so \ref{mk-lem:av_regular_map_is_hom} shows that \(\beta\) is a
  homomorphism.  Thus \(\alpha\) is \(\beta\) followed by translation by \(b\).
\end{proof}
